%! Function Fundamentals: Notation, Domain, and Range — Unit 2, Lesson 2.1 — Homework (Answer Key)
\documentclass[10pt]{article}
\usepackage{atda-article}
\usepackage{atda-key}   % pulls in -boxes; do NOT also load -boxes
\newcommand{\TallMath}[1]{$\displaystyle #1\rule[-1.4em]{0pt}{3.2em}$}
\setlength{\workrowsep}{8pt}

\begin{document}

\pageheader{Unit 2, Lesson 2.1}{Homework --- Answer Key}
% NAMESTRIP: no \namedateperiod here — mirrors the blank. Teacher-only prose goes
% in the lesson plan as \begin{teachernote}[Homework], never in a key.

\vspace{0.15in}

\boxguard[24]
\begin{notesbox}{Practice --- Notation Both Ways, and Writing a Set Down}
\small
Show your work. Items 3 and 4 are the same translation run in opposite directions.

\begin{enumerate}[leftmargin=1.6em, itemsep=6pt, topsep=4pt]

\item Let $f(x)=3x+7$. Find $f(-4)$, then find the value of $x$ for which $f(x)=1$.
\begin{work}
  f(-4) &= 3(-4)+7 \\
        &= -5 \\
   3x+7 &= 1 \\
     3x &= -6 \\
      x &= -2
\end{work}

\item Here is a function $g$ given by a table.
\smallskip

\renewcommand{\arraystretch}{1.25}
\begin{tabular}{@{}|c|c|c|c|c|c|c|c|@{}}
\hline
$x$ & $-3$ & $-2$ & $-1$ & $0$ & $1$ & $2$ & $3$ \\ \hline
$g(x)$ & $5$ & $0$ & $-3$ & $-4$ & $-3$ & $0$ & $5$ \\ \hline
\end{tabular}
\smallskip

Find $g(-1)$. Then find \emph{every} value of $x$ for which $g(x)=0$, and say why there is more
than one.

\ansline{$g(-1)=-3$. \quad $g(x)=0$ at $x=-2$ and at $x=2$.}\\
\ansline{Two inputs may share an output; only the reverse is forbidden.}\\

\item Rewrite each in \textbf{interval} notation: \quad (a) $-3 \le x < 4$ \quad (b) $x > 7$
\quad (c) $x \le 0$

\ansline{(a) $[-3,4)$ \qquad (b) $(7,\infty)$ \qquad (c) $(-\infty,0]$}\\

\item Rewrite each as an \textbf{inequality}: \quad (a) $[2,9]$ \quad (b) $(-\infty, 5]$
\quad (c) $(0, \infty)$

\ansline{(a) $2 \le x \le 9$ \qquad (b) $x \le 5$ \qquad (c) $x > 0$}\\

\item For each function below, state the domain, and name which of the three reasons restricts it
--- \emph{the context}, \emph{the graph}, or \emph{the algebra}. If nothing restricts it, say so.
\begin{enumerate}[label=(\alph*), leftmargin=1.6em, itemsep=3pt, topsep=3pt]
  \item $h(x)=\dfrac{x}{x-8}$
  \item $A(s)=s^2$, the area in square inches of a square tile whose side is $s$ inches
  \item $f(x)=2x-9$, with no situation attached
\end{enumerate}

\ansline{(a) $\{x \mid x \ne 8\}$ --- the algebra: a denominator may not be zero.}\\
\ansline{(b) $s>0$, or $(0,\infty)$ --- the context: a tile has a positive side length.}\\
\ansline{(c) All real numbers, $(-\infty,\infty)$ --- nothing restricts it.}\\

\end{enumerate}
\end{notesbox}

\vspace{0.10in}

\boxguard[30]
\begin{scenariobox}[In Context --- The Highway Arch]{royal}
\small
A road passes under a stone arch. The graph gives $y$, the height of the arch above the road in
feet, at a point $x$ feet to the side of the center line. Negative $x$ is to the left of center.

\begin{center}
\begin{tikzpicture}
\begin{axis}[
    width=10.0cm, height=4.6cm,
    xlabel={$x$ (feet from the center line)}, ylabel={$y$ (height, feet)},
    xmin=-11, xmax=11, ymin=0, ymax=22,
    xtick={-10,-8,-6,-4,-2,0,2,4,6,8,10}, ytick={0,5,10,15,20},
    grid=both, grid style={linegray!45}, axis lines=left,
    tick label style={font=\scriptsize}, label style={font=\scriptsize}]
  \addplot[royal, very thick, domain=-10:10, samples=80]{20-x^2/5};
\end{axis}
\end{tikzpicture}
\end{center}

\begin{enumerate}[leftmargin=1.6em, itemsep=6pt, topsep=2pt, start=6]

\item Read $y(5)$ off the graph. Write it as a point, then as a sentence about the arch.

\ansline{$y(5)=15$, so $(5,15)$ is on the graph.}\\
\ansline{Five feet right of the center line, the arch is 15 feet above the road.}\\

\item Now the other direction: find \emph{every} $x$ for which $y(x)=15$. Say what those two
answers mean for the shape of the arch.

\ansline{$x=-5$ and $x=5$.}\\
\ansline{The arch is symmetric about the center line, so both sides reach 15 ft together.}\\

\item State the domain and the range in this situation, with units, in inequality notation
\emph{and} in interval notation.

\ansline{Domain: $-10 \le x \le 10$ feet, or $[-10,10]$.}\\
\ansline{Range: $0 \le y \le 20$ feet, or $[0,20]$.}\\

\item A moving truck is $12$ feet tall and $10$ feet wide. Driven down the center of the road, its
outer edges are at $x=-5$ and $x=5$. Does it clear the arch? Justify with a value you read off the
graph, not with a guess.

\ansline{Yes. At $x=\pm 5$ --- the truck's widest points --- the arch is 15 feet up.}\\
\ansline{The truck is 12 feet tall, so it clears by 3 feet; everywhere between is higher still.}\\

\end{enumerate}
\end{scenariobox}

\vspace{0.10in}

\boxguard[24]
\begin{scenariobox}[A First Look Ahead]{goldacc}
\small
\begin{enumerate}[leftmargin=1.6em, itemsep=6pt, topsep=2pt, start=10]

\item Two functions, one built from the other: $f(x)=x^2$ and $g(x)=x^2+3$.
\begin{enumerate}[label=(\alph*), leftmargin=1.6em, itemsep=5pt, topsep=4pt]
  \item Fill in the missing row.
        \smallskip

        \renewcommand{\arraystretch}{1.4}
        \begin{tabular}{@{}|c|c|c|c|c|c|@{}}
        \hline
        $x$ & $-2$ & $-1$ & $0$ & $1$ & $2$ \\ \hline
        $f(x)=x^2$ & $4$ & $1$ & $0$ & $1$ & $4$ \\ \hline
        $g(x)=x^2+3$ & \ans{$7$} & \ans{$4$} & \ans{$3$} & \ans{$4$} &
            \ans{$7$} \\ \hline
        \end{tabular}
        \smallskip

  \item Compare the two rows. What happened to \emph{every} output?

\ansline{Every output went up by exactly 3.}\\

  \item Any real number may be squared, so both functions have domain $(-\infty, \infty)$. Write
        the \textbf{range} of each in interval notation. Which one changed --- the domain or the
        range?

\ansline{Range of $f$: $[0,\infty)$. \quad Range of $g$: $[3,\infty)$.}\\
\ansline{Only the range changed; both functions still accept every real number.}\\

  \item Without drawing anything, describe how the graph of $g$ sits compared to the graph of $f$.
        Lesson 2.2 has a name for what you just described.

\ansline{The same parabola, lifted 3 units straight up.}\\
\end{enumerate}

\end{enumerate}
\end{scenariobox}

\vspace{0.10in}

\boxguard[22]
\begin{extensionbox}
\small
\textbf{The booster club, one lesson later.} In Lesson 1.7 the club's average cost per T-shirt was
$A(x)=\dfrac{6x+240}{x}$, where $x$ is the number of shirts ordered. The printer will not take an
order above $200$ shirts.

\smallskip
\renewcommand{\arraystretch}{1.4}
\begin{tabular}{@{}|c|c|c|c|c|@{}}
\hline
$x$ shirts & $10$ & $20$ & $40$ & $60$ \\ \hline
$A(x)$ dollars & \ans{$30$} & \ans{$18$} & \ans{$12$} & \ans{$10$} \\ \hline
\end{tabular}
\smallskip

\begin{enumerate}[label=(\alph*), leftmargin=1.6em, itemsep=5pt, topsep=4pt]
  \item Write the domain of $A$ in this situation in set notation.

\ansline{$\{1, 2, 3, \dots, 200\}$ --- the whole numbers from 1 through 200.}\\

  \item Explain why the domain cannot be written as $(0, 200]$.

\ansline{That interval contains $12.5$, and the club cannot order half a shirt.}\\

  \item In Lesson 1.7 you called $x=0$ an \emph{excluded value} of the expression. Here it is also
        something the \emph{story} explains. Say both things in one sentence.

\ansline{Dividing by $x$ fails at $x=0$, and with no shirts ordered there is no cost per shirt.}\\
\end{enumerate}
\end{extensionbox}

\vspace{0.10in}

\begin{remindbox}
\small
\textbf{Next: Lesson 2.2 --- Parent Functions and Transformations} (\texttt{AFDA.AF.1a--b}).
Item 10 is the whole preview: $g(x)=x^2+3$ is $f(x)=x^2$ with every output pushed up by $3$, which
lifts the range and leaves the domain alone. Lesson 2.2 names that move --- and the three others
like it --- and reads them straight off a graph.
\end{remindbox}

\end{document}
